\section{Dualité de lecture HMT : dualité \texorpdfstring{\(T\)}{T}, bloc \texorpdfstring{\(243\)}{243}, code ternaire de Golay et routes de masse}
\label{sec:sucesor83-f083-dualidad-t}

La dualité \(T\) des cordes se réalise comme projection continue d'une
dualité HMT antérieure : l'échange entre la valeur visible et la mémoire de
frontière. Tout entier admet la lecture
\[
n=9Q+r,
\]
où \(r\) est la valeur visible de la roue nonadique et \(Q\) conserve la
mémoire du franchissement de bord. La même structure apparaît dans l'APP relevée,
dans la décomposition \(1215=54+9\cdot129\), dans l'unité dimensionnelle
stratifiée \(1000=729+243+27+1\) et dans la lecture de la masse comme route avec
avancée visible et enroulement contre-angulaire. Cette dualité relie le bloc
\(243\), la boule ternaire parfaite \(V_3(11,2)=243\), le code ternaire
poinçonné de Golay \([11,6,5]_3\), les canaux de
\(\pi,e,\varphi,\alpha\), la charnière \(744\), la dimension
\(11=6+5=3+8\) et l'involution
\[
(r,Q;R)\longmapsto(Q,r;R^{-1}).
\]
Dans le langage des cordes, la valeur visible est publiée comme mode de vibration ou
impulsion, tandis que la mémoire de frontière est publiée comme enroulement.
En mécanique dimensionnelle, ces deux modes apparaissent comme
\(n_AA_{\mathrm{rad}}\) et \((s_C/6)C^*_{\mathrm{rad}}\) : avancée visible et
mémoire contre-angulaire.

HMT ne prend pas le continu comme objet premier. Elle part d'une cellule discrète
avec mémoire de frontière. La droite réelle, Moonshine et la physique dimensionnelle sont
des lectures différentes de cette même cellule. Par conséquent, la chaîne
\[
\text{APP relevée}
\longrightarrow\text{valeur et mémoire}
\longrightarrow\text{dualité }T
\longrightarrow243/\text{Golay}
\longrightarrow\text{routes de masse}
\]
n'importe pas la dualité de la théorie des cordes : elle montre que sa structure
existe dans HMT avant la géométrisation continue.

\subsection{Reste visible et quotient de frontière}

Pour tout entier positif \(n\), on définit
\[
r_9(n)=1+((n-1)\bmod 9),
\qquad
Q_9(n)=\frac{n-r_9(n)}{9}.
\]
Alors
\[
n=9Q_9(n)+r_9(n).
\]
La lecture ordinaire conserve le reste. La lecture enrichie conserve la
paire complète
\[
n\longmapsto(r,Q).
\]
Cette paire est la structure minimale d'une lecture compacte : une composante est
publiée à l'intérieur de la roue et l'autre conserve le nombre de fois où sa
frontière a été franchie.

\begin{principle}[Dualité entre valeur visible et mémoire de frontière]
La cellule HMT admet deux lectures conjuguées. La première publie \(r\) comme
phase, valeur ou vibration ; la seconde publie \(Q\) comme mémoire,
enroulement ou nombre de franchissements de frontière.
\end{principle}

\subsection{Involution de dualité T}

Soit \(R>0\). La fonctionnelle
\[
E_R(r,Q)=\frac{r^2}{R^2}+Q^2R^2
\]
satisfait
\[
E_R(r,Q)=E_{1/R}(Q,r),
\]
car
\[
E_{1/R}(Q,r)
=\frac{Q^2}{(1/R)^2}+r^2(1/R)^2
=Q^2R^2+\frac{r^2}{R^2}.
\]
Par conséquent,
\[
(r,Q;R)\longleftrightarrow(Q,r;R^{-1}).
\]
L'expression continue des cordes contient des termes de la forme
\[
\frac{n^2}{R^2}+\frac{w^2R^2}{\alpha'^2}.
\]
HMT apporte avant cette réalisation continue le reste visible, la mémoire
de frontière et l'inversion d'échelle. La dualité \(T\) est la projection
géométrique de cette involution de lecture.

\subsection{APP relevée : modes visibles et modes de mémoire}

Pour la somme et le produit sur le même support, on écrit
\[
i+j=9Q^+_{ij}+R^+_{ij},
\qquad
ij=9Q^\times_{ij}+R^\times_{ij}.
\]
L'APP relevée est
\[
\mathcal A_9^\sharp=(R^+,Q^+;R^\times,Q^\times),
\]
où \(R^+,R^\times\) sont les modes visibles et
\(Q^+,Q^\times\) les modes de mémoire. L'écart complet est
\[
ij-(i+j)
=(R^\times_{ij}-R^+_{ij})
+9(Q^\times_{ij}-Q^+_{ij}).
\]
En sommant sur \(1\leq i,j\leq9\),
\[
\sum R^+=405,
\quad
\sum R^\times=459,
\quad
\sum Q^+=45,
\quad
\sum Q^\times=174.
\]
On en obtient
\[
\Delta_{\mathrm{visible}}=459-405=54,
\qquad
\Delta_{\mathrm{memoria}}=174-45=129,
\]
et
\[
\sum ij-\sum(i+j)=2025-810=1215=54+9\cdot129.
\]
La normalisation par la roue produit
\[
\frac{1215}{9}=135=6+129=120+15.
\]
Ainsi, \(54\) est le défaut visible, \(129\) le défaut de mémoire,
\(135\) le défaut total normalisé, \(120\) le canal octadique constitutif
et \(15\) la paire interne de l'hexade.

\subsection{Complétion stratifiée de l'unité}

L'Extrême et Moyenne Raison Holographique possède la complétion cubique
\[
10^3=9^3+3\cdot9^2+3\cdot9+1,
\]
c’est-à-dire,
\[
1000=729+243+27+1.
\]
Les quatre termes représentent respectivement le volume discret de
dimension \(3\), la frontière de codimension \(1\), la frontière de
codimension \(2\) et le retour de codimension \(3\). La couronne active se
raffine en
\[
270=243+27=3\cdot9^2+3\cdot9.
\]
L'unité se décompose ainsi en intérieur, frontière active stratifiée et
retour. Le passage \(9\to10\) ouvre une couche dimensionnelle complète, et la frontière
d'un niveau devient un canal de propagation du niveau suivant.

\subsection{Le bloc \texorpdfstring{\(243\)}{243} comme frontière et boule ternaire parfaite}

Le même bloc s'obtient comme boule de Hamming ternaire de rayon \(2\) et de
longueur \(11\) :
\[
V_3(11,2)
=\sum_{k=0}^{2}\binom{11}{k}2^k
=1+22+220
=243.
\]
Le code ternaire parfait poinçonné de Golay a pour paramètres
\[
G_{11}=[11,6,5]_3,
\qquad
|G_{11}|=3^6=729,
\]
et satisfait
\[
729\cdot243=3^{11}.
\]
Par conséquent, \(G_{11}\) pave \(\mathbb F_3^{11}\) au moyen de boules ternaires de
rayon \(2\). Le nombre \(243\) est simultanément
\[
3\cdot9^2,
\qquad
V_3(11,2),
\qquad
3^5,
\qquad
1+22+220.
\]
Cette identité unit la frontière dimensionnelle, la correction parfaite, la dimension
\(11\), l'incidence de Witt--Golay et la lecture structurelle de la théorie M.

\subsection{Décomposition interne du bloc \texorpdfstring{\(243\)}{243}}

La boule admet la décomposition
\[
243=1+22+90+120+10.
\]
Ici, \(1\) est le centre et le retour ; \(22=2\cdot11\), la première coquille de
Hamming et la lecture électronique du bloc ; \(90=6\binom62\), le canal hexadique ;
\(120=8\binom62\), le canal octadique ; et \(10=9+1\), la clôture de la
monodromie \(9\to1\). Ainsi, le retour, la composante \(22\), les canaux \(90/120\) et la
monodromie sont contenus dans une même boule parfaite.

\subsection{Canaux de constantes dérivés du bloc \texorpdfstring{\(243\)}{243}}

Avec \(B=243\), on obtient les canaux
\[
C_\alpha=3B=729,
\qquad
C_e=B+27+1=271,
\]
\[
C_\pi=B+72=315,
\qquad
C_\varphi=B-81-1=161.
\]
Le canal de \(\alpha\) est le bloc triple \(243\) ; celui de \(e\), le bloc
\(243\) avec la frontière \(27\) et le retour ; celui de \(\pi\), le bloc avec le
canal \(72\) ; et celui de \(\varphi\), le bloc après le retrait du
sous-carré \(81\) et du retour.
La charnière se reconstruit en
\[
C_\pi+C_e+C_\varphi-3
=315+271+161-3
=744
=729+15.
\]
La trace réduite des trois caractères combine donc le bloc triple
\(243\) avec la paire interne de l'hexade.

\subsection{Dimension \texorpdfstring{\(11\)}{11} et lectures compatibles}

Le code ternaire parfait poinçonné produit
\[
11=6+5,
\]
où \(6\) est la dimension de \(G_{11}\) et \(5\) la profondeur de
correction \(3^5=243\). La même dimension se décompose en
\[
11=3+8,
\]
où \(3\) correspond à TRIT et \(8\) à l'octade interne. Dans l'ensemble,
\[
11=6+5=3+8,
\qquad
12=11+1.
\]
La lecture de dimension \(11\) coordonne ainsi code et correction, TRIT et
octade, et l'identité \(729\cdot243=3^{11}\).

\subsection{Réalisation dimensionnelle de la dualité}

L'action locale d'une route s'exprime par
\[
E_0(P)=n_A(P)A_{\mathrm{rad}}
+\frac{s_C(P)}{6}C^*_{\mathrm{rad}}.
\]
Le premier terme est le mode d'avancée ou de vibration visible ; le second, le
mode de frontière ou d'enroulement contre-angulaire. La masse locale est
\[
m_{\mathrm{HMT}}(P)=m_e\exp(E_0(P)).
\]
La tour raffinée ajoute la mémoire de frontière :
\[
E(P)=E_0(P)+k(P)\Delta_4
+\nu_{120}(P)s_{120}
+\nu_{270}(P)\zeta_{270},
\]
où \(\Delta_4\) est la torsion minimale, \(s_{120}\) le sélecteur
hexade--octade et
\[
\zeta_{270}=135\Delta_4^2
\]
la mémoire quadratique de la couronne. La dualité abstraite de route est
\[
E_{A,C}(n,w)=nA_{\mathrm{rad}}+wC^*_{\mathrm{rad}},
\qquad
E_{A,C}(n,w)=E_{C,A}(w,n).
\]
Dans le secteur physique, \(A\) et \(C^*\) sont polarisés, et le registre généalogique sélectionne
la réalisation. La masse est l'exponentielle d'une route de clôture avec avancée
visible, enroulement de frontière et torsion.

\subsection{Compatibilité entre feuillets et clôture du registre généalogique}

Dans une corde fermée, la compatibilité entre les modes gauche et droit
s'exprime schématiquement par
\[
N_L-N_R=nw.
\]
Dans la lecture HMT, le membre de gauche mesure la différence d'excitation
entre feuillets, tandis que celui de droite mesure le couplage entre la phase visible et
la mémoire de frontière. L'ajustement des niveaux constitue ainsi une clôture du
registre généalogique bidirectionnel. Dans les lectures \(CP\), \(T\) et \(CPT\), une asymétrie
d'orientation est compensée par l'inversion bidirectionnelle pour clore le
registre.

\subsection{Dualité T, modularité et fonction J}

Dans le tore complexe,
\[
q=e^{2\pi i\tau},
\]
et la transformation modulaire
\[
\tau\longmapsto-\frac1\tau
\]
échange ses cycles. La dualité \(T\) échange les grands et les petits
rayons ; les deux réalisations publient l'échange entre le cycle visible
et le cycle de frontière. La fonction modulaire \(J\), l'action de
Sommerfeld et la dualité \(T\) constituent des clôtures distinctes d'une même
structure première de frontière avec mémoire.

\subsection{Conséquences dimensionnelles pour le photon, l'électron, la masse et la statistique}

Le photon observable est associé au mode central du TPK :
\[
L_{\mathrm{TPK}}=3(I-P_0),
\qquad
\operatorname{Spec}(L_{\mathrm{TPK}})=\{0,3,3\}.
\]
Le mode zéro est une propagation nulle étendue sans mémoire compactifiée
propre. L'électron apparaît comme clôture compacte des modes nuls
conjugués,
\[
e^-=\text{fermeture compacte de }P_++P_-,
\]
et constitue l'enroulement physique minimal avec mémoire stable. La masse
mesure le coût de stabilisation de cette mémoire compacte. L'exclusion de
Pauli s'interprète comme saturation du registre généalogique : deux fermions identiques n'
occupent pas la même frontière parce que leur clôture d'enroulement est déjà saturée.
Les bosons publient des modes partageables de mémoire déployée, tandis que les
fermions publient une mémoire compactifiée.

\subsection{Identités vérifiables de la construction}

La construction satisfait conjointement les identités
\begin{enumerate}
\item \(1000=729+243+27+1\);
\item \(V_3(11,2)=1+22+220=243\);
\item \(729\cdot243=3^{11}\);
\item \(243=1+22+90+120+10\);
\item \(C_\alpha,C_e,C_\pi,C_\varphi=729,271,315,161\);
\item \(315+271+161-3=744\);
\item \(1215=54+9\cdot129\);
\item \(E_R(r,Q)=E_{1/R}(Q,r)\);
\item \(11=6+5=3+8\), \(12=11+1\), \(26=2+24\) y \(10=2+8\).
\end{enumerate}

La dualité \(T\) est donc reconnue comme la réalisation géométrique de l'
échange entre reste et quotient. L'enroulement n'est pas introduit
d'abord comme propriété extérieure d'une corde compacte, mais comme mémoire de
frontière de la décomposition \(n=9Q+r\). Le bloc \(243\) réunit cette
dualité avec Golay ternaire, la dimension \(11\), les canaux \(90/120\), la génération
de constantes, les routes de masse et la théorie M. Dans sa forme la plus générale, HMT lit la
masse, la corde et la dimension comme des publications d'une même frontière avec
mémoire.
